\chapter{Food Labels}

\section*{Definition and Overview}
Food labels are the information panels on packaged foods that tell consumers what is in the product, how much of each nutrient it contains, and how it should be stored and used. Learning to read food labels is one of the most useful skills for healthy eating, because it allows people to compare products, limit salt, sugar, and saturated fat, identify allergens, and choose foods that match their health goals. Australian food labels include the ingredient list, the nutrition information panel, the health star rating, and allergen declarations, and the Food Standards Code requires these to be accurate and truthful. This chapter explains how to read each part of a food label and use it to make healthier choices.

\section*{Epidemiology}
Food labels are on virtually all packaged foods in Australia, but their use varies. Surveys show that most shoppers read labels at least sometimes, most often the ingredient list and use-by date, while fewer use the nutrition information panel to compare fat, sugar, and salt. The Health Star Rating system, introduced to simplify choices, now appears on a large share of eligible products. Misreading labels is common: many people do not know that sugar has many different names, that "light" does not always mean low in energy, and that the per-serve values can hide high nutrient content. Better label literacy is associated with healthier food choices.

\section*{Aetiology and Pathophysiology}
The elements of a label each have a purpose.
\begin{itemize}
    \item \textbf{Product name and claims:} descriptions and marketing claims, which must not be misleading but require checking against the panel
    \item \textbf{Ingredient list:} all ingredients, listed in descending order of weight, so the first ingredients are the most abundant
    \item \textbf{Nutrition information panel:} energy, protein, total and saturated fat, carbohydrate, sugars, sodium, and fibre, shown per serve and per 100 g
    \item \textbf{Per 100 g comparison:} the standard way to compare products, because serve sizes differ between brands
    \item \textbf{Health Star Rating:} a 0.5 to 5 star system based on the product's nutrients, providing an at-a-glance guide
    \item \textbf{Allergen declaration:} a mandatory statement of the major allergens (peanut, tree nuts, milk, egg, wheat, soy, fish, shellfish, and sesame), essential for people with allergies
    \item \textbf{Date marking:} use-by dates for perishable foods and best-before dates for others
    \item \textbf{Storage and preparation:} instructions that affect food safety
    \item \textbf{Nutrient claims:} terms such as "low fat", "high fibre", and "no added sugar", which have defined legal meanings but need checking against the panel
\end{itemize}

\section*{Clinical Features}
Skilled label reading has recognisable features.
\begin{itemize}
    \item Comparing products by the per 100 g values rather than the serve sizes
    \item Checking total and saturated fat, sugars, and sodium in the nutrition panel
    \item Reading the ingredient list to see the main ingredients and hidden sugars
    \item Using the Health Star Rating as a quick guide, and the panel for detail
    \item Checking the allergen statement carefully for food allergies
    \item Recognising that "no added sugar" does not mean no sugar, and "light" does not mean low energy
    \item Understanding that claims such as "natural" and "organic" do not necessarily mean healthier
    \item Using use-by and best-before dates for food safety
\end{itemize}

\section*{Differential Diagnosis}
Common label-reading questions need careful answers.
\begin{itemize}
    \item "Low fat" products, which may be high in sugar
    \item "No added sugar" products, which can still contain natural sugars
    \item Sugars under many names: sucrose, glucose, fructose, maltose, honey, and syrups
    \item Sodium hidden as salt, or in additives
    \item The difference between total fat and saturated fat
    \item Health Star Rating versus the detailed panel, which may differ for whole foods
    \item Allergen statements versus "may contain" statements, which are precautionary
\end{itemize}

\section*{When to Seek Medical Care}
Reading labels is a self-care skill, but people with food allergies, coeliac disease, diabetes, heart disease, or kidney disease should seek dietitian advice on label reading for their specific needs.

\textbf{Call 000 (or local emergency services) for:}
\begin{itemize}
    \item Anaphylaxis from accidental allergen exposure, despite label checking
    \item Food poisoning from spoiled or incorrectly stored food
    \item Any medical emergency
\end{itemize}

\section*{Diagnosis and Investigations}
\begin{itemize}
    \item \textbf{Per 100 g check:} use these numbers to compare fat, sugar, salt, and fibre between brands
    \item \textbf{Ingredient list check:} the first ingredients are the most abundant, and hidden sugars appear under many names
    \item \textbf{Health Star Rating:} use for at-a-glance comparison of similar products
    \item \textbf{Allergen check:} the mandatory allergen statement for people with allergies
    \item \textbf{Portion awareness:} check the serve size and how it compares to what is actually eaten
    \item \textbf{Dietitian review:} for individualised label reading for medical conditions
\end{itemize}

\section*{Management}
\begin{itemize}
    \item \textbf{Compare per 100 g:} aim for products lower in saturated fat, sugar, and sodium, and higher in fibre
    \item \textbf{Check the Health Star Rating:} prefer 3.5 stars or more where available, while checking the panel
    \item \textbf{Read the ingredient list:} choose foods with shorter lists and recognisable ingredients
    \item \textbf{Watch serve sizes:} a product can look healthy per serve while containing a lot per 100 g
    \item \textbf{Look for hidden sugar:} recognise sugar's many names
    \item \textbf{Check allergen statements:} always, for people with allergies
    \item \textbf{Use dates and storage instructions:} to prevent food poisoning and waste
    \item \textbf{Remember whole foods:} unpackaged fruit, vegetables, and unprocessed foods need no label at all
\end{itemize}

\section*{Complications}
\begin{itemize}
    \item Choosing unhealthy products based on marketing claims rather than the panel
    \item Accidental allergen exposure from missed label reading
    \item Excess energy intake from "low fat" products high in sugar
    \item High salt and sugar intake from misjudged "healthy" options
    \item Food waste and illness from ignored use-by dates
\end{itemize}

\section*{Prognosis and Outcomes}
People who read labels well make measurably healthier choices, with lower intake of saturated fat, sugar, and salt, and better management of allergies, diabetes, and heart disease. Label reading is a skill that improves with practice, and the Health Star Rating has made it easier for the general population. Combined with a diet built around whole foods, which need no labels, accurate label reading supports long-term health and the prevention of chronic disease.

\section*{Prevention and Self-Care}
\begin{itemize}
    \item Always compare per 100 g values between similar products
    \item Use the Health Star Rating as a quick guide
    \item Learn the common names for sugar and salt
    \item Check the allergen statement for every product if you have an allergy
    \item Note the serve size and how it relates to what you actually eat
    \item Prefer whole foods, which have no label
    \item Ask a dietitian for label-reading help for medical conditions
    \item Check use-by dates and storage instructions for food safety
\end{itemize}

\section*{Sources and Further Reading}
The following reputable sources provide further information for healthcare students and patients seeking reliable, current advice about food labels.
\begin{itemize}
    \item Food Standards Australia New Zealand. \textit{Food labelling.} https://www.foodstandards.gov.au/
    \item Healthdirect Australia. \textit{How to read food labels.} https://www.healthdirect.gov.au/
    \item Health Star Rating. \textit{Understanding the Health Star Rating.} http://www.healthstarrating.gov.au/
    \item Dietitians Australia. \textit{Reading food labels.} https://dietitiansaustralia.org.au/
    \item NHS. \textit{How to read food labels.} https://www.nhs.uk/live-well/eat-well/
    \item World Health Organization. \textit{Food labelling.} https://www.who.int/
\end{itemize}
